\documentclass[twoside]{article}

\usepackage[preprint]{aistats2027}
\usepackage[round,authoryear]{natbib}
\usepackage{amsmath,amssymb,amsthm}
\usepackage{mathtools}
\usepackage{microtype}
\usepackage{booktabs}
\usepackage{url}

\renewcommand{\bibname}{References}
\renewcommand{\bibsection}{\subsubsection*{\bibname}}
\bibliographystyle{apalike}

\newtheorem{theorem}{Theorem}
\newtheorem{proposition}[theorem]{Proposition}
\newtheorem{corollary}[theorem]{Corollary}
\newtheorem{lemma}[theorem]{Lemma}
\newtheorem{remark}[theorem]{Remark}
\newcommand{\E}{\mathbb E}
\newcommand{\R}{\mathbb R}
\newcommand{\cF}{\mathcal F}
\newcommand{\RR}{\mathsf{RR}}
\newcommand{\SGD}{\mathsf{SGD}}
\newcommand{\one}{\mathbf 1}
\newcommand{\asympup}{\mathrel{\lesssim}}
\newcommand{\asympdown}{\mathrel{\gtrsim}}

\begin{document}
\runningauthor{}

\twocolumn[
\aistatstitle{Minimax Random Reshuffling in Every Convex Regime}

\aistatsauthor{Pahan Dewasurendra}

\aistatsaddress{Johns Hopkins University\\[4pt] July 30, 2026}
]

\begin{abstract}
Random reshuffling is the standard implementation of stochastic gradient descent on finite sums, but its finite-time theory was incomplete. A recent upper bound shows that, for every reasonable constant step size and every number of epochs, averaging all within-epoch iterates incurs \[ \frac{D^2}{\eta nK} +\min\{\eta,\eta^2nL\}\sigma_\star^2. \] It was open whether this interpolation between stochastic gradient descent and reshuffling is sharp.

We prove a matching lower bound after the natural smoothness cap. Consequently, the minimax error of optimally tuned constant-step random reshuffling is \[
\begin{aligned}
 \Theta\!\Bigg(\min\!\Bigg\{&
 LD^2,\frac{LD^2}{nK}\\[-0.5ex]
 &+\min\!\left\{\frac{\sigma_\star D}{\sqrt{nK}},
 \left(\frac{L\sigma_\star^2D^4}{nK^2}\right)^{1/3}\right\}
 \Bigg\}\Bigg).
\end{aligned}
\] The result holds for every $K$ and every $n\ge32$. It closes both restrictions in the previous convex lower bound, which needed $\eta=O(1/(nL))$ and many epochs and concerned epoch endpoints. We also match the separate average and maximum component smoothness parameters whenever the effective active count $n\bar L/\widehat L$ is at least a constant.

The proof uses an asymmetric random-permutation bridge. Its new ingredient is to set the active curvature to $\ell=\min\{L,1/(161\eta n)\}$. The same small-effective-step bridge then produces $\eta^2nL\sigma_\star^2$ for small steps and $\eta\sigma_\star^2$ for large steps. A common contraction inequality controls every within-epoch iterate, while an orthogonal deterministic coordinate covers the finite-horizon transient. The hard instances are three-dimensional, component convex, and have exact smoothness and variance parameters.
\end{abstract}

\section{Introduction}
\label{sec:intro}

For a finite sum $f=n^{-1}\sum_{i=1}^n f_i$, practical stochastic gradient descent usually processes a random permutation of the components in each epoch. This random reshuffling strategy is easy to implement and often outperforms independent sampling. Its analysis is harder because a component gradient is not conditionally unbiased after the preceding elements of a permutation have been observed.

A long line of work established faster asymptotic rates under small steps or large epoch counts \citep{nagaraj2019sgd,haochen2019random, mishchenko2020random,ahn2020sgd,nguyen2021unified}. Matching lower bounds revealed a genuine reshuffling advantage, first for quadratics and later for smooth components \citep{safran2020good,rajput2020closing,safran2021many,cha2023tighter}. These results did not settle the smooth convex problem at practical constant steps or at finite epoch counts.

\citet{liu2026dominates} recently removed both restrictions on the upper-bound side. Let $D=\lVert x_0-x_\star\rVert$, let $L$ bound each component smoothness, and define
\begin{equation}
 \sigma_\star^2=\frac1n\sum_{i=1}^n
 \lVert\nabla f_i(x_\star)\rVert^2.
 \label{eq:sigma}
\end{equation}
For any $\eta\le 1/(6L)$, their theorem bounds the objective error at the average of all within-epoch iterates by
\begin{equation}
 O\!\left(
 \frac{D^2}{\eta nK}
 +\min\{\eta,\eta^2nL\}\sigma_\star^2
 \right).
 \label{eq:liu-fixed}
\end{equation}
The first noise term is the familiar $\SGD$ scale. The second is the faster reshuffling scale. The minimum proves that $\RR$ is never worse than $\SGD$ for any reasonable step. The paper explicitly asks for a complete lower bound under arbitrary reasonable steps and finite $K$.

We answer this question for optimally tuned constant-step $\RR$. Up to universal constants, its minimax error is
\begin{equation}
 \min\!\left\{LD^2,\frac{LD^2}{nK}
 +\min\!\left\{
 \frac{\sigma_\star D}{\sqrt{nK}},
 \left(\frac{L\sigma_\star^2D^4}{nK^2}\right)^{1/3}
 \right\}\right\}.
 \label{eq:rate-intro}
\end{equation}
The outer cap is unavoidable. As $\eta$ tends to zero, the output tends to the initial point, whose objective gap is at most $LD^2/2$.

\paragraph{Why prior constructions stop at small steps.}
The sharp convex lower bound of \citet{cha2023tighter} uses a piecewise quadratic whose curvature is larger on the negative half-line. A balanced random permutation makes the iterates trace a discrete bridge. The curvature asymmetry converts the bridge magnitude into a one-sided bias. Their movement estimates require $\eta nL$ to be a small constant. Their theorem also needs enough epochs for the bias to persist at epoch endpoints.

Our key reduction is to separate the class smoothness $L$ from the curvature that creates the bridge. For a tested step $\eta$, set
\begin{equation}
 \ell=\min\left\{L,\frac{1}{161\eta n}\right\}.
 \label{eq:ell}
\end{equation}
The bridge always sees a small effective step. Its error scale is \[ \ell\eta^2n\sigma_\star^2 =\min\{\eta^2nL,\eta/161\}\sigma_\star^2, \] which is precisely the two-branch noise term in \eqref{eq:liu-fixed}, up to a constant. Thus one construction covers both small and large steps.

Two further ideas handle the output and finite horizon. First, we strengthen the asymmetric bridge from epoch endpoints to every within-epoch iterate. A pointwise comparison on negative starts and an expected-gradient inequality on positive starts admit one common contraction coefficient. Second, an orthogonal slow quadratic coordinate supplies the transient whenever the bridge has not accumulated. A dormant $L$-quadratic coordinate makes every component's minimal smoothness constant exactly $L$ without changing the dynamics.

\paragraph{Contributions.}
\begin{itemize}
\item We prove the fixed-step lower bound
\[ \Omega\!\left( \min\!\left\{LD^2,\frac{D^2}{\eta nK} +\min\{\eta,\eta^2nL\}\sigma_\star^2\right\}\right) \] for every $\eta\le1/(6L)$ and every finite $K$.
\item Optimizing the step gives the matching minimax rate
\eqref{eq:rate-intro}. This closes the small-step and large-epoch restrictions of the previous convex lower bound.
\item The hard instance controls the exact all-within-iterate average used by
the new upper bound. It has dimension three, convex components, exact $\sigma_\star$, and exact component smoothness $L$.
\item An active-subset construction extends the rate to heterogeneous
smoothness, matching both $\bar L$ and $\widehat L$ in the upper bound when $n\bar L/\widehat L\ge34$.
\end{itemize}

\section{Setting and Upper Benchmark}
\label{sec:setting}

For $d\ge3$, let $f_i:\R^d\to\R$ be differentiable, convex, and $L$-smooth. Write $f=n^{-1}\sum_i f_i$ and assume that $f$ has a minimizer $x_\star$. Random reshuffling starts from $x_{1,0}$ and, independently in each epoch $k$, draws a uniform permutation $\pi_k$ of $[n]$ and updates
\begin{equation}
 x_{k,i}=x_{k,i-1}-\eta\nabla f_{\pi_k(i)}(x_{k,i-1}),
 \qquad i\in[n].
 \label{eq:rr}
\end{equation}
Set $x_{k+1,0}=x_{k,n}$. The output is
\begin{equation}
 \bar x=\frac{1}{nK}\sum_{k=1}^K\sum_{i=0}^{n-1}x_{k,i}.
 \label{eq:output}
\end{equation}
This is exactly the constant-step specialization of the weighted output of \citet{liu2026dominates}.

Let $\cF_n(L,D,\sigma)$ contain all such finite sums, over arbitrary $d\ge3$, paired with an initialization satisfying
\begin{equation}
 \lVert x_{1,0}-x_\star\rVert\le D,
 \qquad \sigma_\star\le\sigma,
 \label{eq:class}
\end{equation}
where $\sigma_\star$ is defined in \eqref{eq:sigma}. Define the optimally tuned constant-step minimax risk
\begin{equation}
\begin{aligned}
 \mathfrak M_{n,K}(L,D,\sigma)
 =\inf_{0<\eta\le1/(6L)}
 \sup_{\substack{(f,x_{1,0})\\
 \in\cF_n(L,D,\sigma)}}\\[-0.5ex]
 \hspace{2.8em}\E[f(\bar x)-f(x_\star)].
\end{aligned}
 \label{eq:minimax}
\end{equation}

The upper side of our result follows from the recent theorem and a uniform small-step limit.

\begin{proposition}[Upper benchmark \citealp{liu2026dominates}]
\label{prop:upper}
For every $n,K\ge1$,
\begin{align}
 \mathfrak M_{n,K}(L,D,\sigma)
 \asympup
 \min\!\left\{LD^2,\frac{LD^2}{nK}\right.\quad\notag\\[-1ex]
 \left.{}+\min\!\left\{
 \frac{\sigma D}{\sqrt{nK}},
 \left(\frac{L\sigma^2D^4}{nK^2}\right)^{1/3}
 \right\}\right\}.
 \label{eq:upper}
\end{align}
\end{proposition}

The $LD^2$ term follows from a uniform iterate bound as $\eta$ tends to zero. The proof is in the supplement. The remaining terms are Corollary 1 of \citet{liu2026dominates} with equal component smoothness constants. We retain the cap because it is needed in low-noise or short-horizon regimes.

\section{Main Results}
\label{sec:results}

Our fixed-step theorem matches the dependence in \eqref{eq:liu-fixed} whenever that expression is informative relative to smoothness.

\begin{theorem}[All-regime fixed-step lower bound]
\label{thm:fixed}
There is a universal $c>0$ such that, for every $n\ge32$, $K\ge1$, $L,D,\sigma>0$, and $0<\eta\le1/(6L)$, there is an instance in $\cF_n(L,D,\sigma)$ with $\sigma_\star=\sigma$, $\lVert x_{1,0}-x_\star\rVert=D$, and component smoothness exactly $L$ for which
\begin{align}
 \E[f(\bar x)-f(x_\star)]
 \ge c\min\!\left\{LD^2,\frac{D^2}{\eta nK}\right.\quad\notag\\[-1ex]
 \left.{}+\min\{\eta,\eta^2nL\}\sigma^2\right\}.
 \label{eq:fixed-lower}
\end{align}
The instance has dimension three. It may depend on the tested step $\eta$, as is standard for a minimax lower bound.
\end{theorem}

The restriction $n\ge32$ only simplifies numerical constants. The supplement also gives the dummy-component reduction for odd $n$.

Combining Theorem~\ref{thm:fixed} with Proposition~\ref{prop:upper} gives the complete optimized characterization.

\begin{corollary}[Minimax rate]
\label{cor:minimax}
For every $n\ge32$ and $K\ge1$,
\begin{align}
 \mathfrak M_{n,K}(L,D,\sigma)
 \asymp
 \min\!\left\{LD^2,\frac{LD^2}{nK}\right.\quad\notag\\[-1ex]
 \left.{}+\min\!\left\{
 \frac{\sigma D}{\sqrt{nK}},
 \left(\frac{L\sigma^2D^4}{nK^2}\right)^{1/3}
 \right\}\right\}.
 \label{eq:minimax-rate}
\end{align}
The constants in $\asymp$ are universal.
\end{corollary}

The stochastic branches meet at
\begin{equation}
 K=\frac{L^2D^2n}{\sigma^2}.
 \label{eq:transition}
\end{equation}
Below this threshold, the optimal step satisfies $\eta nL\gtrsim1$ and the $\SGD$ branch is minimax. Above it, $\eta nL\lesssim1$ and reshuffling gives the cubic branch. No assumption that $K$ exceeds the threshold is made.

\subsection{Separate average and maximum smoothness}
\label{sec:heterogeneous}

The upper theorem of \citet{liu2026dominates} uses \[ \bar L=\frac1n\sum_{i=1}^nL_i, \qquad \widehat L=\max_iL_i \] separately. The maximum restricts the step, while the average controls the small-step noise. The equal-smoothness result is not exploiting an ambiguity in these parameters. We can match the full two-parameter rate.

Let $L_i$ denote the minimal gradient-Lipschitz constant of $f_i$. Let $\cF_n(\bar L,\widehat L,D,\sigma)$ be the analogue of the class in Section~\ref{sec:setting}, with these constants satisfying the display above. Define
\begin{equation}
\begin{aligned}
 &\mathfrak M_{n,K}(\bar L,\widehat L,D,\sigma)\\
 &\quad=\inf_{0<\eta\le1/(6\widehat L)}
 \sup_{\cF_n(\bar L,\widehat L,D,\sigma)}
 \E[f(\bar x)-f(x_\star)].
\end{aligned}
 \label{eq:minimax-hetero}
\end{equation}

\begin{theorem}[Heterogeneous minimax rate]
\label{thm:heterogeneous}
If $n\bar L/\widehat L\ge34$, then
\begin{align}
 \mathfrak M_{n,K}(\bar L,\widehat L,D,\sigma)
 \asymp
 \min\!\left\{\bar LD^2,\frac{\widehat LD^2}{nK}\right.\quad\notag\\[-1ex]
 \left.{}+\min\!\left\{
 \frac{\sigma D}{\sqrt{nK}},
 \left(\frac{\bar L\sigma^2D^4}{nK^2}\right)^{1/3}
 \right\}\right\}.
 \label{eq:minimax-hetero-rate}
\end{align}
The constants are universal.
\end{theorem}

Here is the construction. Choose an even $m\asymp n\bar L/\widehat L$. Only these $m$ components act on the hard $u,v$ coordinates. The others are no-ops there. Give active components smoothness $\widehat L$ in the dormant coordinate, and choose the inactive dormant curvature so that the average is exactly $\bar L$. Put balanced active gradients of magnitude $\sigma\sqrt{n/m}$ at the optimum. Their full $n$-component variance is exactly $\sigma^2$.

Deleting no-ops from a uniform permutation leaves a uniform permutation of the active components. More subtly, the output still averages all $n$ states. Conditional on the active order, inactive positions form a uniform interleaving. If $z_0,\ldots,z_m$ are the active states, the expected coefficient on each $z_r$ for $r<m$ is \[ \frac{n+1}{n(m+1)}. \] The endpoint coefficient is nonnegative. Since the bridge gives $\E z_r\ge0$, the full all-iterate mean is at least one half of the active all-iterate mean.

Apply Lemma~\ref{lem:bridge} with \[ \ell=\min\{\widehat L,1/(161\eta m)\}. \] The full objective contains $m h/n$, so its stochastic scale is \[ \frac mn\ell\eta^2 \left(\sigma\sqrt{\frac nm}\right)^2m =\ell\eta^2m\sigma^2 \asymp\min\{\eta^2n\bar L,\eta\}\sigma^2. \] Using the same active subset for the slow coordinate gives $\min\{\bar LD^2,D^2/(\eta nK)\}$. The finite-horizon case split and scalar optimization then prove Theorem~\ref{thm:heterogeneous}. Full details, including the choice of $m$ and exact interleaving coefficients, are in the supplement.

\section{The Rescaled Asymmetric Bridge}
\label{sec:bridge}

We isolate the stochastic part of the proof. Let $n$ be even and choose balanced $b_i\in\{-\sigma,+\sigma\}$. For an active curvature $\ell>0$, set
\begin{equation}
 h(v)=\frac{\ell}{2}(v_-)^2
 +\frac{\ell}{2\cdot2415}(v_+)^2,
 \label{eq:h}
\end{equation}
where $v_+=\max\{v,0\}$ and $v_-=\max\{-v,0\}$. The components are $h_i(v)=h(v)+b_iv$. Their average is $h$, and $v_\star=0$.

The following lemma strengthens the endpoint mechanism of \citet{rajput2020closing} and \citet{cha2023tighter} to the output \eqref{eq:output}.

\begin{lemma}[All-iterate asymmetric bridge]
\label{lem:bridge}
Suppose $n\ge32$ is even, $\eta\ell n\le1/161$, and $v_{1,0}=0$. Run $K$ epochs of $\RR$ on the components $h_i$, and let $\bar v$ average all pre-update iterates. There is a universal $c_b>0$ such that
\begin{equation}
 \E\bar v\ge c_b\eta\sigma\sqrt n
 \min\{1,K\eta\ell n\}.
 \label{eq:bridge-mean}
\end{equation}
Consequently,
\begin{equation}
 \E h(\bar v)\ge c_b'
 \ell\eta^2n\sigma^2
 \min\{1,K\eta\ell n\}^2.
 \label{eq:bridge-risk}
\end{equation}
\end{lemma}

\paragraph{Proof idea.}
Write $Y_k=v_{k,0}$, $M_k=\E Y_k$, and $\alpha=\eta\ell n$. The bridge estimates of \citet{cha2023tighter} imply that, conditional on a fixed nonnegative start, the first half of an epoch has expected iterate
\begin{equation}
 \E[v_{k,i}\mid Y_k=y]
 \ge \left(1-\frac23\eta\ell i\right)y
 +\frac{\eta^2\ell\sigma}{480}
 \sum_{j=0}^{i-1}\sqrt j.
 \label{eq:first-half}
\end{equation}
For the second half, their two expected-gradient estimates give
\begin{equation}
 \E[v_{k,i}\mid Y_k=y]
 \ge(1-0.36\alpha)y
 +\frac{\eta^2\ell\sigma n^{3/2}}{18000}.
 \label{eq:second-half}
\end{equation}
The numerical constants are deliberately loose.

Negative starts require care because replacing $y$ by $M_k$ in these inequalities is not automatic. Couple the asymmetric process to the symmetric $\ell$-quadratic process under the same permutation. The former is pointwise larger after every update. Hence, for $y<0$,
\begin{equation}
 \E[v_{k,i}\mid Y_k=y]\ge(1-\eta\ell)^iy.
 \label{eq:negative}
\end{equation}
For the first half, $(1-\eta\ell)^i\le1-(2/3)\eta\ell i$. In the second half, $(1-\eta\ell)^i\le1-0.36\alpha$. Multiplication by $y<0$ reverses these coefficient inequalities. Thus positive and negative starts share the same coefficient before averaging over $Y_k$.

A sign-flip coupling gives $\Pr(Y_k\ge0)\ge1/2$ in every epoch. The endpoint version of \eqref{eq:second-half} then yields
\begin{equation}
 M_{k+1}\ge(1-0.36\alpha)M_k
 +\frac{\alpha\eta\sigma\sqrt n}{36000}.
 \label{eq:epoch-recursion}
\end{equation}
Starting from $M_1=0$, this recursion keeps $M_k$ nonnegative and gives
\begin{equation}
 \frac1K\sum_{k=1}^KM_k
 \asympdown \eta\sigma\sqrt n\min\{1,(K-1)\alpha\}.
 \label{eq:mean-starts}
\end{equation}
At least one third of the within-epoch iterates also receive the direct positive term in \eqref{eq:second-half}. Averaging all iterates gives \eqref{eq:bridge-mean}, including $K=1$. Since $h(v)\ge\ell v^2/(2\cdot2415)$, Jensen's inequality gives \eqref{eq:bridge-risk}. The supplement proves every displayed inequality and tracks integer endpoints.

The curvature choice \eqref{eq:ell} is now decisive. It ensures the lemma's small-effective-step condition for every $\eta$, while
\begin{equation}
 \ell\eta^2n
 =\min\{\eta^2nL,\eta/161\}
 \asymp\min\{\eta^2nL,\eta\}.
 \label{eq:two-branches}
\end{equation}
This turns one local bridge inequality into both stochastic regimes.

\section{Proof of the Lower Bound}
\label{sec:proof}

We prove Theorem~\ref{thm:fixed} for even $n$. The odd case is in the supplement. Let $T=nK$ and define
\begin{equation}
 \lambda=\min\left\{L,\frac{1}{2\eta T}\right\},
 \qquad \ell=\min\left\{L,\frac{1}{161\eta n}\right\}.
 \label{eq:curvatures}
\end{equation}
On $x=(u,v,w)\in\R^3$, set
\begin{equation}
 f_i(u,v,w)=\frac\lambda2u^2+h(v)+b_iv+\frac L2w^2,
 \label{eq:hard}
\end{equation}
with $h$ from \eqref{eq:h}, and initialize at $(D,0,0)$. The average objective is $\lambda u^2/2+h(v)+Lw^2/2$ and has minimizer zero. Every component is convex. Its minimal smoothness constant is exactly $L$ because of the dormant $w$ coordinate. Also, the initial distance is exactly $D$ and $\sigma_\star=\sigma$. The $w$ coordinate remains zero.

The $u$ iterates satisfy $u_t=(1-\eta\lambda)^tD$. Since $\eta\lambda T\le1/2$, every pre-update iterate is at least $D/2$. Therefore
\begin{equation}
 \frac\lambda2\bar u^2
 \asympdown \min\left\{LD^2,\frac{D^2}{\eta T}\right\}.
 \label{eq:transient}
\end{equation}

Lemma~\ref{lem:bridge} and \eqref{eq:two-branches} control the $v$ coordinate. If $\eta L T\ge1/161$, then \[ K\eta\ell n=K\min\{\eta nL,1/161\}\ge1/161. \] The attenuation in \eqref{eq:bridge-risk} is then a universal constant, so
\begin{equation}
 \E h(\bar v)
 \asympdown\min\{\eta,\eta^2nL\}\sigma^2.
 \label{eq:stochastic-full}
\end{equation}
Combining \eqref{eq:transient} and \eqref{eq:stochastic-full} gives \eqref{eq:fixed-lower}. If $\eta LT<1/161$, then $\lambda=L$ and \eqref{eq:transient} alone is $\Omega(LD^2)$, which again gives \eqref{eq:fixed-lower}. This proves the theorem.

For completeness, optimizing the lower bound is elementary. Put \[ \Phi(\eta)=\frac{D^2}{\eta nK} +\sigma^2\min\{\eta,\eta^2nL\}. \] On $\eta\ge1/(nL)$, balancing the first term with $\eta\sigma^2$ gives $\eta=D/(\sigma\sqrt{nK})$ and value $\sigma D/\sqrt{nK}$. On $\eta\le1/(nL)$, balancing with $\eta^2nL\sigma^2$ gives
\begin{equation}
 \eta=\left(\frac{D^2}{n^2KL\sigma^2}\right)^{1/3}
 \label{eq:eta-cubic}
\end{equation}
and the cubic term in \eqref{eq:minimax-rate}. The two candidate steps lie in their stated regions on opposite sides of the threshold \eqref{eq:transition}. Clipping at $1/(6L)$ contributes $LD^2/(nK)$. Finally, the outer minimum with $LD^2$ commutes with the infimum over $\eta$. This proves Corollary~\ref{cor:minimax}.

\section{Relation to Prior Lower Bounds}
\label{sec:related}

\citet{safran2020good} analyze final epoch iterates of strongly convex quadratics and establish sharp weighted-permutation variance estimates. \citet{safran2021many} add condition-number dependence and show that reshuffling need not improve on independent sampling before many epochs. These quadratic constructions are symmetric and target final iterates.

\citet{rajput2020closing} introduced the asymmetric piecewise quadratic for a strongly convex endpoint lower bound. \citet{cha2023tighter} sharpened the mechanism, gave lower bounds for arbitrary weighted averages of epoch endpoints, and extended the cubic rate to convex objectives. Their convex theorem assumes $\eta\le1/(161nL)$ and a large epoch count. Our bridge uses their local permutation estimates with attribution. The common-coefficient argument, control of every within-epoch iterate, curvature rescaling, dormant smoothness coordinate, and finite-horizon direct sum are new.

Recent papers improve upper bounds through primal-dual structure, heterogeneous smoothness, or last-iterate analysis \citep{cai2024tighter,liu2024last}. They do not supply the lower bound considered here. The later small-epoch study of \citet{kim2025incremental} concerns deterministic incremental gradients and lists random-reshuffling lower bounds as open. Most decisively, \citet{liu2026dominates} proves the exact upper benchmark used here and states that a complete lower characterization for reasonable steps and finite epochs is missing.

Two adjacent preprints address different questions. \citet{li2025unified} unify nonconvex upper analyses through an order-error assumption. \citet{peng2026resolution} resolve operator-norm comparisons among shuffle schemes for near-identity quadratics. Neither gives an objective-risk lower bound with gradient noise.

\section{Discussion}
\label{sec:discussion}

The rate has a simple interpretation. Random reshuffling has two stochastic time scales. When the stable memory $1/(\eta L)$ is shorter than an epoch, the hard instance behaves at the $\SGD$ noise scale. When it is longer, balancing a random-permutation bridge against curvature gives the cubic reshuffling scale. The active-curvature reduction makes this interpolation visible in one hard family rather than through unrelated case constructions.

Our result is information theoretic and concerns the algorithm exactly as stated. The hard instance may depend on its constant step. We do not claim a lower bound against arbitrary time-varying schedules, momentum, variance reduction, or adaptive permutation design. Last-iterate risks can differ from the all-iterate output studied here. Extending the characterization to nonconvex components or fixed componentwise smoothness profiles remains open.

The proof uses no high-dimensional embedding and no oracle hiding. Three coordinates suffice: one for deterministic transient, one for permutation bias, and one dormant coordinate that certifies the exact class smoothness. This low-dimensional structure suggests that the transition in \eqref{eq:minimax-rate} reflects the dynamics of reshuffling itself rather than statistical complexity of the function class.

\bibliography{references}

\clearpage
\end{document}
